\section{Study 1a: The cross-domain reliability landscape}
% =====================================================================
\subsection{Method (summary)}
Study 1a systematically reviewed \Nstudies{} studies across \Ndomains{} domains comparing LLM-as-a-judge to human
annotation. Identification triangulated OpenAlex (\OAretrieved{} records, \OAunique{} unique), Semantic Scholar, and
a 59-query web search; screening retained on-topic empirical comparisons; \Nnumeric{} studies reported an
extractable numeric agreement value (Figure~\ref{fig:prisma1}; the full coded corpus is \corpustabref).
\textbf{Coding reliability:} because the verdict rubric embeds judgement, we validated it with an independent
re-coding of a 25\% random sample ($n{=}\IRRsoneN$): the reliability verdict reproduced at Krippendorff
$\alpha=\IRRsoneA$ (ordinal) with \IRRsoneWone\% of cases within one level (Appendix~\ref{app:irr}). The re-coders are
LLM-assisted, like the primary pipeline, so this measures rubric \emph{reproducibility} rather than human
inter-rater agreement, which we flag as a limitation. \textbf{Citation integrity:} every one of the
$\Nstudies{+}\STN$ references was independently web-verified for existence against arXiv/venue/DOI records;
$270/272$ were confirmed (260 at high confidence) and the $2$ that could not be verified were removed, so the
corpus contains no unverifiable citations.

\begin{figure}[t]\centering
\includegraphics[width=0.96\linewidth]{fig01_prisma_flow}
\caption{Study 1a PRISMA~2020 flow (two-arm identification: OpenAlex database + web/citation).}
\label{fig:prisma1}
\end{figure}

\subsection{Findings: a green/amber/red landscape}
The modal verdict is \emph{mixed/task-dependent}; only \Pctreliable\% of studies sit in the reliable band and
\Pctproblematic\% in the problematic band (mean reliability \Meanrel/5). Reliability is governed by three properties:
\textbf{objectivity}, \textbf{tractable verification}, and a \textbf{low human-agreement floor}.
\textbf{Green ($\ge$3.5; \Ngreen{} domains):} machine translation (system level), text-to-image alignment,
verifiable instruction-following, QA/factuality, information extraction, general-NLG pairwise preference, dialogue,
and legal scoring: here LLM judges frequently match or exceed the human ceiling (e.g., MT-Bench 85\% vs.\ 81\%
human-human \citep{zheng2023mtbench}; GEMBA-MQM 96.5\% \citep{kocmi2023gembamqm}; SAFE wins 76\% of disagreements
\citep{wei2024safe}). \textbf{Amber (\Namber{} domains):} summarization, RAG, education, social-science annotation,
medicine (structure-dependent: $\kappa\!\approx\!0.82$-$0.93$ structured vs.\ $0.47$-$0.71$ complex reasoning),
code, and IR/relevance (strong but contested \citep{clarke2024cantreplace}). \textbf{Red (\Nred{} domains):}
low-resource multilingual evaluation and creative writing (near-zero expert correlation
\citep{chakrabarty2024artartifice}). Figures~\ref{fig:dom1}-\ref{fig:land1} give the domain ordering and the
domain$\times$task landscape; the decision map (Figure~\ref{fig:dec1}) positions domains by reliability against
stakes.

\begin{figure}[t]\centering
\begin{subfigure}{0.5\linewidth}\includegraphics[width=\linewidth]{fig05_reliability_by_domain}\caption{}\label{fig:dom1}\end{subfigure}\hfill
\begin{subfigure}{0.48\linewidth}\includegraphics[width=\linewidth]{fig06_reliability_landscape}\caption{}\label{fig:land1}\end{subfigure}
\caption{Study 1a. (a) Mean reliability by domain (green/amber/red). (b) Reliability landscape: domain$\times$judging
task; verification/classification columns are greener than open-ended scoring.}
\end{figure}

\subsection{Failure modes and the decision map}
Recurring biases such as verbosity/length, position, self-preference, style-over-substance, preference leakage,
non-determinism, and the \emph{expertise gap} (judges converge to crowd, not expert, consensus
\citep{bavaresco2024judgebench}) all modulate reliability (Figure~\ref{fig:bias1}). Agreement should be benchmarked
against the human ceiling, which on subjective tasks is itself low (Krippendorff $\alpha<0.35$
\citep{imrit2026iaa}). The Study-1a decision map (Figure~\ref{fig:dec1}) and Table~\ref{tab:decision} translate this into
deployment guidance and motivate both the moving-target test (Study 1b) and the HCC-specific analysis that follows.

\begin{figure}[t]\centering
\begin{subfigure}{0.5\linewidth}\includegraphics[width=\linewidth]{fig11_bias_frequency}\caption{}\label{fig:bias1}\end{subfigure}\hfill
\begin{subfigure}{0.49\linewidth}\includegraphics[width=\linewidth]{fig15_decision_map}\caption{}\label{fig:dec1}\end{subfigure}
\caption{Study 1a. (a) Catalogue of reported failure modes/biases. (b) Deployment decision map (reliability$\times$
stakes; bubble area $\propto$ studies).}
\end{figure}

\begin{table}[tbp]\centering\small
\caption{Domain-level decision table. Mean reliability score (1 = unreliable $\dots$ 5 = validated), modal verdict,
and traffic-light signal (\textsc{green} $\ge 3.5$; \textsc{amber} 2.5-3.5; \textsc{red} $<2.5$). Sorted by mean
reliability. Small-$n$ domains should be read together with the family-level aggregates.}
\label{tab:decision}
\begin{tabular}{@{}l r c l c@{}}
\toprule
Domain / task & $n$ & Mean rel. & Modal verdict & Signal \\
\midrule
Machine Translation & 3 & 4.33 & Promising & \cellcolor{g!55}\textsc{green} \\
Text-to-Image & 2 & 4.00 & Promising & \cellcolor{g!55}\textsc{green} \\
Dialogue & 1 & 4.00 & Promising & \cellcolor{g!55}\textsc{green} \\
Instruction Following & 3 & 4.00 & Validated & \cellcolor{g!55}\textsc{green} \\
QA \& Factuality & 5 & 3.80 & Promising & \cellcolor{g!55}\textsc{green} \\
Information Extraction (NER/NLI) & 9 & 3.78 & Promising & \cellcolor{g!55}\textsc{green} \\
General NLG & 21 & 3.71 & Promising & \cellcolor{g!55}\textsc{green} \\
Legal & 2 & 3.50 & Promising & \cellcolor{g!55}\textsc{green} \\
Education \& Grading & 17 & 3.35 & Mixed & \cellcolor{a!60}\textsc{amber} \\
RAG & 3 & 3.33 & Mixed & \cellcolor{a!60}\textsc{amber} \\
Summarization & 8 & 3.25 & Promising & \cellcolor{a!60}\textsc{amber} \\
Fact-Checking & 4 & 3.25 & Promising & \cellcolor{a!60}\textsc{amber} \\
Sentiment \& Emotion & 4 & 3.25 & Mixed & \cellcolor{a!60}\textsc{amber} \\
IR \& Relevance & 9 & 3.22 & Promising & \cellcolor{a!60}\textsc{amber} \\
Social-Science Annotation & 15 & 3.20 & Mixed & \cellcolor{a!60}\textsc{amber} \\
Peer Review \& Science & 6 & 3.17 & Mixed & \cellcolor{a!60}\textsc{amber} \\
Multimodal (Vision-Language) & 5 & 3.00 & Mixed & \cellcolor{a!60}\textsc{amber} \\
Qualitative Coding & 6 & 3.00 & Mixed & \cellcolor{a!60}\textsc{amber} \\
Argument Quality & 1 & 3.00 & Mixed & \cellcolor{a!60}\textsc{amber} \\
Finance & 3 & 3.00 & Mixed & \cellcolor{a!60}\textsc{amber} \\
Medicine \& Clinical & 16 & 2.88 & Mixed & \cellcolor{a!60}\textsc{amber} \\
Methodology \& Bias & 38 & 2.87 & Mixed & \cellcolor{a!60}\textsc{amber} \\
Reward Modeling & 5 & 2.80 & Mixed & \cellcolor{a!60}\textsc{amber} \\
Code \& SW Eng & 7 & 2.71 & Mixed & \cellcolor{a!60}\textsc{amber} \\
Reasoning \& Math & 3 & 2.67 & Mixed & \cellcolor{a!60}\textsc{amber} \\
Content Moderation \& Toxicity & 5 & 2.60 & Mixed & \cellcolor{a!60}\textsc{amber} \\
Mental Health & 4 & 2.50 & Mixed & \cellcolor{a!60}\textsc{amber} \\
Creative Writing & 6 & 2.33 & Mixed & \cellcolor{r!55}\textsc{red} \\
Multilingual & 4 & 2.25 & Caution & \cellcolor{r!55}\textsc{red} \\
\bottomrule
\end{tabular}
\end{table}

% =====================================================================
